\chapter[Modern Cosmology and the Cosmological Constant]{Modern Cosmology and the\\Cosmological Constant\label{chap_MCCC}}
\epigraph{It is of great advantage to the student of any subject to read the original memoirs on that subject, for science is always most completely assimilated when it is in the nascent state.}{---James Clerk Maxwell, \textit{A Treatise on Electricity and Magnetism}}
\section{Early Modern Cosmology\label{sec_EMC}}
The era of modern cosmology, in our continued search for an accurate understanding of the large-scale structure and evolution of the Universe, began only a century ago, when advances were made in physical theory due to Albert Einstein's general relativity theory \cite{Einstein1915,Einstein1952}, and, shortly thereafter, to the observations of cosmic expansion. This expansion was first evidenced by redshift measurements of spiral nebulae, after the task of measuring their radial velocities was initiated in 1912 by Vesto Slipher \cite{Slipher1913}. Within a decade, the compiled list of these measurements, being predominantly recessional, strongly evidenced the expansion of the Universe \cite{Eddington1923}. The expansion was eventually confirmed by Edwin Hubble, who, after using observations of Cepheid variable stars to show that the spiral nebulae were distant galaxies, in 1926 \cite{Hubble1926a,Hubble1926b}, determined, in 1929, a linear velocity-distance relation among the nearby galaxies \cite{Hubble1929}. 

Prior to Hubble's discovery, few models had been proposed as potential descriptors of the Universal geometry. In particular, the static, homogeneous, finite geometry proposed in 1917 by the father of general relativity theory \cite{Einstein1917,Einstein1952}, required a modification of the original Einstein equation,
\begin{equation}
G_{\mu\nu}\equiv{R_{\mu\nu}}-\frac{1}{2}Rg_{\mu\nu}=\kappa{T_{\mu\nu}},\label{Einstein_incomplete}
\end{equation}
by the addition of a cosmological term,
\begin{equation}
G_{\mu\nu}+\Lambda{g_{\mu\nu}}=\kappa{T_{\mu\nu}},\label{Einstein_complete}
\end{equation}
in order to balance otherwise contracting physical distances between fundamental observers in the matter-filled Universe; i.e., `necessary \ldots for the purpose of making possible a quasi-static distribution of matter, as required by the fact of the small velocity of the stars.' The constant $\Lambda$ in Eq.~(\ref{Einstein_complete}) is called the \textit{cosmological constant}, and the $\Lambda$-term is the only term up to second order which may be added to Einstein's original theory, Eq.~(\ref{Einstein_incomplete}), without violating general covariance, and thus local energy conservation.

An alternative to Einstein's static solution for the Universe was found by Willem de~Sitter in the same year that Einstein published his model, in a successful attempt to discover a solution with a `preferred' elliptical space, as opposed to the spherical space in Einstein's solution \cite{deSitter1917}; however, as Arthur Eddington noted a few years later, this discrepancy between the two spatial geometries is really due to an arbitrarily introduced division of space and time, and that Einstein's \textit{spacetime} geometry is cylindrical, with time uncurved, while de~Sitter's is formally equivalent to a 4-sphere, say of radius $R$, as its line-element may be written,
\begin{equation} 
ds^2=R^2[d\omega^2+\sin^2\omega\{d\zeta^2+\sin^2\zeta(d\theta^2+\sin^2\theta{d\phi^2})\}],\label{deSitter_sphere}
\end{equation}
where $\omega$, $\zeta$, $\theta$, and $\phi$ are four angular coordinates \cite{Eddington1923}. 

De~Sitter's solution contains no matter, and had therefore been a blow to Einstein who had hoped Mach's principle, along with the principle of relativity expressed by general covariance and the principle of equivalence, would be one of three cornerstones of the correct theory of gravitation; i.e.~he hoped the geometry would be completely determined by the stress-energy of matter \cite{Pais1982}. Abraham Pais \cite{Pais1982} writes, 
\begin{quotation}
Einstein never said so explicitly, but it seems reasonable to assume that he had in mind that the correct equations should have no solutions at all in the absence of matter. However, right after his paper appeared, de~Sitter did find a solution \ldots with $\rho=0$. \ldots In 1918 he looked for ways to rule out the de~Sitter solution, but soon realised there is nothing wrong with it.
\end{quotation}

De~Sitter's seminal paper of 1917 \cite{deSitter1917} contains the first attempt at a relativistic interpretation of the observed velocities of the spiral nebulae, previously evidenced by their redshifts. Here, he analysed the static solutions of Einstein's equations for which the matter was supposed a pressureless perfect fluid, denoted \textit{dust} in modern parlance, with constant stress-energy $T_{00}=\rho$, else $T_{\mu\nu}=0$, and found that if the density was negligible (in fact, he set $\rho=0$) the equations, Eq.~(\ref{Einstein_complete}), are satisfied by
\begin{equation}
ds^2=-\cos^2\chi{c^2}dt^2+R^2d\chi^2+R^2\sin^2\chi(d\theta^2+\sin^2\theta{d\phi^2}),\label{deSitter}
\end{equation} 
with $\Lambda=3/R^2$. De~Sitter compared his solution to the Minkowski and Einstein geometries, which also satisfy Eq.~(\ref{Einstein_complete}) with the given stress-energy, and found that if his geometry would be the correct representation of World-structure we should observe redshifts in the spectra of distant objects as `the frequency of light-vibrations diminishes with increasing distance from the origin of co-ordinates' due to $g_{00}=\cos^2\chi$, so that the time for a cycle of vibrations would be proportional to $\sec\chi$. He therefore proposed that this redshift might be measured in the spiral nebulae, which were not yet known as distant galaxies similar to the Milky Way, but, due to their measured radial velocities, were nevertheless thought to be the most distant objects known. In 1917, de~Sitter was able to quote only three reliable radial velocity measurements of the nebulae, which gave $2:1$ odds in favour of his prediction.

Later, Eddington further analysed de~Sitter's solution and found a change of basis to the static, spherically symmetric line-element \cite{Eddington1923},
\begin{equation}
ds^2=-\left(1-\frac{\Lambda}{3}r^2\right)c^2dt^2+\frac{dr^2}{1-\frac{\Lambda}{3}r^2}+r^2(d\theta^2+\sin^2\theta{d\phi^2}).\label{deSitter_statical}
\end{equation}
This solution had in fact been found in 1918 by Friedrich Kottler \cite{Kottler1918}, who generalised Karl Schwarzschild's solution for a static, spherically symmetric point mass \cite{Schwarzschild1916,Schwarzschild2003} to allow for the possibility of non-zero $\Lambda$. Kottler's solution,
\begin{equation}
ds^2=-\left(1-\frac{2M}{r}-\frac{\Lambda}{3}r^2\right)c^2dt^2+\frac{dr^2}{1-\frac{2M}{r}-\frac{\Lambda}{3}r^2}+r^2(d\theta^2+\sin^2\theta{d\phi^2}),\label{SdS_statical}
\end{equation}
where $M$ is called the \textit{mass parameter}, reduces to a representation either of Schwarzschild's or de~Sitter's geometry when $\Lambda=0$ or $M=0$, respectively, as Eddington had finally shown in 1923 \cite{Eddington1923}; thus, the geometry Kottler had first described generally is known as the Schwarzschild-de~Sitter (SdS) geometry, and we refer to Eq.~(\ref{SdS_statical}) as the statical SdS metric. It will be central to our analysis in Chapter~\ref{chap_ANC}, through which it will be shown to describe a universe that would theoretically appear as ours does.

Now, from Eq.~(\ref{deSitter_statical}) Eddington found that the redshift de~Sitter had predicted as a phenomenon of his static geometry was in fact due to a repulsion brought in by $\Lambda$ \cite{Eddington1923}; viz., he showed that particles could not remain at constant physical distances from each other in the geometry described by Eq.~(\ref{deSitter_statical}), but would recede according to
\begin{equation}
\frac{d^2r}{ds^2}=\frac{\Lambda}{3}r,
\end{equation}
so that while every particle would claim to remain at rest (at $r=0$) in de~Sitter space, each would say that all others are exponentially receding. Then, due to the implicit symmetry of the geometry---i.e.~from Eq.~(\ref{deSitter_sphere})---Eddington concluded that all observers in de~Sitter's geometry would observe this same phenomenon, and independently arrived at an expanding solution by 1923. This he substantiated with an updated list of recession velocities of the spiral nebulae provided by Slipher, which gave $36:41$ odds in favour of the expansion.

Eddington argued that the singularity de~Sitter had found at $\chi=\pi/2$---a `paradox' which de~Sitter had dismissed, as Eddington put it, because its effects would be relevant only to events `before the beginning or after the end of eternity' \cite{Eddington1923}---could not be physical; it is coordinate-dependent, existing at a radius $r=R=\sqrt{3/\Lambda}$ from every observer in the statical solution, Eq.~(\ref{deSitter_statical}}), and thus the entire sphere can be covered by intersections of the `lunes' of neighbouring particles, though they could never share information received from outside each other's respective `lune'. 

Eddington's argument sufficiently removed the paradox of the de~Sitter horizon, but his solution still contained the coordinate singularity, and he had even suggested that this singularity could not be removed in a meaningful coordinate basis \cite{Eddington1923}. This would eventually be proven false: the first coordinate transformation which removed the singularity was found by Georges Lema\^{i}tre in 1925, while he studied at MIT \cite{Lemaitre1925b}. Eddington had noticed, in \cite{Eddington1923}, that de~Sitter's original solution, Eq.~(\ref{deSitter}), did not contain the proper time coordinate, and was therefore not a good representation of the geometry, and proposed that the static line-element, Eq.~(\ref{deSitter_statical}), used the correct basis. However, what Lema\^{i}tre noticed was that in de~Sitter's solution constant spatial coordinates are not geodesics---de~Sitter's frame is not \textit{comoving}. Instead, he noted that de~Sitter's system has a preferred location at $\chi=0$. This is equally true of the origin in the static line-element proposed by Eddington. What Lema\^{i}tre discovered instead, was that a reference frame which treats all points in space on equal footing---a truly comoving system of geodesics in de~Sitter space---is an exponentially expanding Galilean system, in which the coordinate singularity occurs at $t=\pm\infty$. In modern notation, this solution is written
\begin{equation}
ds^2=-dt^2+e^{2Ht}[dr^2+r^2(d\theta^2+\sin^2\theta{d}\varphi^2)],\label{deSitter_comoving}
\end{equation}
where $H\equiv1/R=\sqrt{\Lambda/3}$, and we see that the geometry is both homogeneous and isotropic, as in Einstein's model. Lema\^{i}tre would conclude \cite{Lemaitre1925b},
\begin{quotation}
Our treatment evidences [the] non-statical character of de~Sitter's world which gives a possible interpretation of the mean receding motion of spiral nebulae.

\ldots[However,] we are led back to the euclidean space and to the impossibility of filling up an infinite space with matter which cannot but be finite. De~Sitter's solution has to be abandoned, not because it is non-static, but because it does not give a finite space without introducing an impossible boundary.
\end{quotation}

This represents an important step in the progression of Lema\^{i}tre's understanding, which unfortunately went largely unnoticed by cosmologists at the time---his paper having been largely unknown. Due to Einstein's and Eddington's influence, Lema\^{i}tre, among others, had believed very strongly in the concept of a closed, unbounded spherical Universe, which he forever retained in his cosmological models. In the publication of his static model in 1917 \cite{Einstein1917}, Einstein had noticed that in order to describe the Universe it would be incorrect to assume, as he did in his calculation of the perihelion precession of Mercury, that space would be asymptotically flat. And the ingenuity with which Einstein solved this problem, by introducing his spherical universe, was noted later by Eddington \cite{Eddington1933}:
\begin{quotation}
I think Einstein showed his greatness in the simple and drastic way in which he disposed of difficulties at infinity. He abolished infinity. He slightly altered his equations so as to make space at great distances bend round until it closed up. So that, if in Einstein's space you keep going right on in one direction you do not get to infinity; you find yourself back at your starting-point again. Since there was no longer an infinity, there could be no difficulties at infinity. \textsc{q.e.d.}
\end{quotation}
Eddington was convinced already by 1923 \cite{Eddington1923} that $\Lambda$ is a Universal constant which gives the radius of curvature of the vacuum at every point, as $\sqrt{3/\Lambda}$, and that `The length of a specified material structure bears a constant ratio to the radius of curvature of the world at the place and in the direction in which it lies'; thus, he associated the fact that material objects are measured to have the same lengths, regardless of location, with the positive constancy of $\Lambda$. He summarised his argument,
\begin{quotation}
From this point of view it is inevitable that the constant $\Lambda$ cannot be zero; so that empty space has a finite radius of curvature relative to familiar standards. An electron could never decide how large it ought to be unless there existed some length independent of itself for it to compare itself with.
\end{quotation}
Lema\^{i}tre was well aware of, and much influenced by both of these theories, and therefore sought a universal model which would account for both, as Einstein's model did, but which would also account for the redshifts of the nebulae, which only an `impossibly bounded' de~Sitter model could do. He therefore sought `an intermediate solution which could combine the advantages of both' the Einstein and de~Sitter universes \cite{Lemaitre1931a}. This, he achieved in 1927 \cite{Lemaitre1927}, in a paper which would unfortunately remain unnoticed until after he brought it to Eddington's attention a few years later, following Hubble's discovery.

But Lema\^{i}tre's solution was in fact a rediscovery of an earlier solution found by Aleksandr Friedman in 1922, which was unknown to most western astronomers and theoreticians until after Hubble's discovery. In two separate papers \cite{Friedmann1922,Friedmann1999a,Friedmann1924,Friedmann1999b}, Friedman published models for isotropic, homogeneous, expanding/contracting matter-filled universes with constant positive and negative curvature, whose dynamical equations (written in modern notation),
\begin{equation}
3\frac{\dot{a}^2+K}{a^2}-\Lambda=\kappa\rho\label{Friedmann}
\end{equation}
\begin{equation}
2\frac{\ddot{a}}{a}+\frac{\dot{a}^2+K}{a^2}-\Lambda=-\kappa{p},\label{Friedmann2}
\end{equation}
where $\kappa=8\pi{G}/c^4$ is the proportionality constant from the Einstein equation, Eq.~(\ref{Einstein_complete}), describe a time-dependent scale factor, $a=a(t)$, typically scaled by requiring the value of the curvature constant, $K$, to be $-1$ or $+1$. Note that Friedman did not explicitly treat the Euclidean case, for which $K=0$; nor did he account for non-zero pressure. The solution Lema\^{i}tre had rediscovered in \cite{Lemaitre1927} was the case $K=+1$, with generalisation to $p\neq0$, in accordance with his understanding of the Universe. He further analysed the astronomical aspects of this solution, providing a description of the receding velocities of the nebulae according to the expansion of space \cite{Krasinski1999b}. 

The pressureless, dust-filled Einstein model may be recovered from Eqs.~(\ref{Friedmann}) and~(\ref{Friedmann2}) by requiring $\ddot{a}=\dot{a}=0$, $\rho=\rho_m$, and $p=p_m=0$, for which we obtain the result,
\begin{equation}
\Lambda_{E}=\frac{\kappa}{2}\rho_m=K/a_{E}^2=+1/a_E^2.\label{a_E}
\end{equation}
Lema\^{i}tre's comoving de~Sitter geometry, Eq. (\ref{deSitter_comoving}), is equivalent to Friedman's solution with $a(t)=e^{Ht}$ and $p=\rho=K=0$. All are special cases of the general form of the Robertson-Walker metrics, Eq.~(\ref{RW}), discussed at the end of this section.

As mentioned above, the works of Friedman and Lema\^{i}tre went effectively unnoticed throughout the 1920s. Einstein had known of Friedman's original article, but at first thought Friedman had made a calculation error. In a note he submitted shortly after its publication \cite{Einstein1922}, he claimed that Friedman's result seemed suspicious (`verd\"{a}chtig'); however, in a retraction he submitted the following year \cite{Einstein1923}, after Friedman had managed to deliver him an explanation \cite{Krasinski1999b}, he called Friedman's result correct and clarifying (`richtig und aufkl\"{a}rend'). In spite of this, Einstein would make no further reference to Friedman's work until after Hubble had sufficiently demonstrated the Universal expansion; indeed, the Einstein and de~Sitter universes remained the two contending geometries of cosmological theories until 1930. For instance, in 1926, following the publication of his famous result for the distance to the Andromeda nebula, Hubble began considering the possibility of a static, homogeneous spatial distribution of galaxies, and from his result estimated the radius of curvature for the Einstein Universe, $a_E=2.7\times10^{10}$~pc \cite{Hubble1926b}. Then, in 1929, he made note of the de~Sitter expansion as the possible cause of the linear redshift-distance relation \cite{Hubble1929}, which was a known first-order approximation of the redshift in de~Sitter's geometry---the relevant calculation of this, by Howard Robertson in 1928 \cite{Robertson1928}, which contains references to the earlier work on this subject, is discussed below.

Then, in 1930 Lema\^{i}tre made known to Eddington his previous result, \cite{Lemaitre1927}. In his famous paper in which he proved Einstein's static universe was unstable \cite{Eddington1930}, Eddington begins:
\begin{quotation}
Working in conjunction with Mr.~G.~C.~McVittie, I began some months ago to examine whether Einstein's spherical universe is stable. Before our investigation was complete we learnt of a paper by Abb\'{e} G.~Lema\^{i}tre which gives a remarkably complete solution of the various questions connected with the Einstein and de~Sitter cosmogonies.
\end{quotation} 
Then, in the same month as Eddington's article had appeared (May), de~Sitter published an article in which Friedman is credited as the original discoverer of the expanding universe solution \cite{deSitter1930}:
\begin{quotation}
A dynamical solution \ldots is given by Dr.~G.~Lema\^{i}tre in a paper published in 1927, which had unfortunately escaped my notice until my attention was called to it by Professor Eddington a few weeks ago.
\end{quotation}
In de~Sitter's citation to Lema\^{i}tre's article \cite{Lemaitre1927}, he adds a note: `A similar solution had previously been given by A.~Friedman \ldots', citing \cite{Friedmann1922}.

Following this exposition of the previously unnoticed works on the expansion paradigm, there was a flurry of articles published which would model the expanding Universe. In March 1931, Lema\^{i}tre published an (updated) English translation of his 1927 article \cite{Lemaitre1931a}, along with a further analysis of his considerations of the beginning of expansion from a dense Einstein universe \cite{Lemaitre1931b}, which elaborated on Eddington's previous analysis \cite{Eddington1930}. In this paper, Lema\^{i}tre generalised the proof of the Jebsen \cite{Jebsen1921,Jebsen2005}-Birkhoff \cite{Birkhoff1923} theorem (see \cite{Deser2005}) to non-zero $\Lambda$, showing that the geometry outside a (nonrotating and uncharged) spherical mass is SdS, even if expanding or contracting. Thus, Lema\^{i}tre proved that the formation of condensations in an Einstein universe would not lead directly to expansion or contraction of the universe \cite{Lemaitre1931b}. Then in May of the same year, Lema\^{i}tre submitted his proposal for `The Beginning of the World from the Point of View of Quantum Theory' \cite{Lemaitre1931c}; from a singularity which Fred Hoyle would later dub the `Big Bang'. Lema\^{i}tre's later work, e.g.~\cite{Lemaitre1933,Lemaitre1997,Lemaitre1934}, would be aimed at describing the present Universe as having long since expanded beyond its Einstein radius, where it had remained for some time in a nearly static state due to only a very slight imbalance in favour of the cosmic repulsion of the $\Lambda$-term over the Newtonian attraction. 

However, in April 1931, Einstein published a paper which would take a markedly different approach. In his now-famous retraction of the $\Lambda$-term \cite{Einstein1931}, he suggested a cycloidal model\footnote{Note that in \cite{Einstein1931} Einstein concentrates only on one period of the cycloid to describe our Universe, and makes no mention of the possibility of a periodic behaviour, as in `big bounce' or `phoenix' theories. The periodic behaviour of this universe was in fact first suggested by Friedman in \cite{Friedmann1922,Friedmann1999a}, although Einstein's analysis in \cite{Einstein1931} had perhaps made this more plain. However, credit for the first serious proposal of an oscillating universe must be given to Richard Tolman, who presented a thermodynamic model of the sort in 1931 \cite{Tolman1931}. In this paper, Tolman writes, `Recently, \ldots a simple model of the universe has been discussed by Einstein \cite{Einstein1931} which exhibits a possibility for a quasi-periodic solution of a type which must now also be considered.' It is likely the ambiguity of this statement which has led some---e.g., \cite{Penrose1979a}---to credit Einstein with the suggestion; for one can easily mis-interpret from what Tolman wrote, that Einstein had made note of the possibility of quasi-periodicity. In fact, Tolman's reference to Einstein's work ends at the citation.} for the Universe---the well-known `closed Einstein-de~Sitter model'---which was essentially a re-statement of Friedman's original solution for $\Lambda=0$, substantiated by Hubble's discovery. A few months later, in July, Otto Heckmann published the first suggestion within the expanding Universe paradigm that $\Lambda$ and the spatial curvature $K$ might also be negative or zero \cite{Heckmann1931}. 

Note, that the solutions studied by Heckmann, on which he elaborated in a second article published the following February \cite{Heckmann1932}---e.g., plotting all solutions for the nine general models in which relativistic and non-relativistic matter do not interact,---had been published by Robertson already in October 1929 \cite{Robertson1929}, well after Hubble had sufficiently demonstrated the linear velocity-distance relation of the galaxies \cite{Hubble1929}. We must therefore dwell briefly on this paper by Robertson, because it is in fact very unclear, from it alone, which stance Robertson had taken on the matter of expansion---although it has been remarked that he had `discarded [the expanding models] in favour of static universes' \cite{Krasinski1999b},---and because there is a direct link to the proposal of one of the most important aspects of all cosmological theories, known as `Weyl's postulate' \cite{Robertson1933,Bondi1961, Rugh2010}.

The successful goal of \cite{Robertson1929}  was to derive, in general, the isotropic and homogeneous geometries. However, Robertson failed to describe the redshift by formulating the evolution of the scale factor in a meaningful way---e.g., by relating his results to Friedman's equations. The paper spends little effort in describing the non-stationary solutions, aside from including them in general formulations. Instead, Robertson always makes special effort to relate his results to a second assumption, that the Universe be stationary. From these assumptions alone, he proved that the only possible stationary cosmologies (with positive curvature) are those of Einstein and de~Sitter. Among these, he shows that only the de~Sitter cosmology can account for the observed redshifts. He seems not to have taken preference to any of the solutions he had derived, but had kept to his goal of an analytical proof that geometries which had previously been known, but had previously only been given as solutions which would satisfy the field equations in a specific way, rather than as special cases of a more general solution, are the only ones that exist.

However, his tendency to give the stationary solutions special treatment, along with the result that Einstein's universe contains no systematic redshift, is suggestive of a preference towards the de~Sitter solution over Friedman's and Einstein's. This is evident in the work he had then been carrying out in the same vein as Hermann Weyl. In 1923, Weyl published an analysis of de~Sitter's geometry---first, in the fifth appendix to his book \textit{Raum, Zeit, Materie} (\textit{Space, Time, Matter}) \cite{Weyl1923a}, which was never translated to English, and then in a supplementary paper \cite{Weyl1923,Weyl2009}; see the editor's note, \cite{Ehlers2009}, to the reprint of \cite{Weyl1923}---which was unique from Eddington's approach in that it described an expanding kinematical system, from which two key advancements were made: the first calculation of redshift in a non-static solution (which, to first order, anticipated the linear redshift-distance relation found later by Hubble \cite{Hubble1929}), and the proposal of a coherency of matter which describes the concept of a cosmic standard time, the importance of which was noted in particular by Robertson \cite{Robertson1929,Robertson1933}, who had referred to it as `Weyl's postulate' in \cite{Robertson1933}, and then later exposed by Hermann Bondi \cite{Bondi1961}, who is credited with actually giving it this name. On this postulate, Weyl wrote that the assumption that two stars would `belong to the same system and [be] causally connected through a common origin, implies that their world lines have the same action domain $\Sigma$.\symbolfootnote[1]{Both the papers by de~Sitter \cite{deSitter1917} and Eddington \cite{Eddington1923} lack this assumption on the ``state of rest'' of stars---by the way the only possible one compatible with the homogeneity of space and time. Without such an assumption nothing can be known about the redshift, of course. [Weyl's footnote.]}' \cite{Weyl2009}. Weyl's model defined the universe as the three-dimensional slice of de~Sitter spacetime that Lema\^{i}tre later described with his line-element, Eq.~(\ref{deSitter_comoving}), which would cover half of the entire `sphere', bounded by the infinitely distant past and infinitely distant future---c.f.~the singularity at $\chi=\pi/2$ in de~Sitter's coordinates, Eq.~(\ref{deSitter}),---of which any observer would see only a `cuneiform sector' \cite{Weyl1930}. This universe was to have begun an infinite time in the past, when the distances between stars would tend to zero, and would expand according to the cosmic repulsion described by $\Lambda$ \cite{Weyl2009}:
\begin{quote}
All stars of our system $\Sigma$ flee from any arbitrary star in radial directions; there is inherent in matter a universal tendency to expand which finds its expression in the ``cosmological term'' of \textsc{Einstein's} law of gravitation \ldots.
\end{quote}

In 1928, Robertson proposed a formally equivalent model to Weyl's \cite{Robertson1928}, in which he operated from comoving coordinates derived independently from Lema\^{i}tre. In a first-order calculation, he derived a linear velocity-distance relation, to which he noted Weyl's result would also reduce. This cosmological theory was tidied up in a paper written by Weyl which was submitted in July 1929 \cite{Weyl1930}, after he had visited both Robertson and Tolman in the United States. He proposed that the Universe began from a common origin at $t=-\infty$ and subsequently expanded as a homogeneous matter distribution of infinitely small density with open space and infinite mass, and updated his derivation of the observed redshift using the comoving coordinates. It was likely this theory which Robertson had in the back of his mind when he wrote \cite{Robertson1929}, which may have been the reason that in it he paid little attention to the non-static solutions. However, the theory never became popular, probably because \cite{Weyl1930} was published in May 1930---the same month as Eddington's and de~Sitter's papers on the expanding Universe, which exposed Lema\^{i}tre's previous work.

Thus, although it seems that Robertson had favoured the de~Sitter solution in \cite{Robertson1929}, this cannot have been because he had hoped it would describe a non-expanding picture of the Universe. Hubble's demonstration of the linear velocity-distance relation of extra-galactic nebulae had been published eight months prior, and had actually been anticipated by Robertson, among others. Instead, there seems to be an implicit assumption in \cite{Robertson1929}, that the true physical cosmology must have a stationary geometrical representation. This bias is not explicitly voiced, but it seems reasonable to speculate that it did exist, and that it was related to one of the cosmological conditions that Tolman had suggested in April that year \cite{Tolman1929a}; viz.~that it be possible to write the line element `in a form which is spherically symmetrical in the spatial variables, symmetrical with respect to past and future time, and static with respect to time.' Tolman's physical motivation for proposing such a condition was not that he had thought the Universe was properly static, but that, respectively, each element of the condition would require a universe which, on the large scale, did not have different properties in different directions, would have large-scale reversible behaviour, and would, by-and-large, be in a steady state \cite{Tolman1929a}. These requirements would be rejected within a year, when it was understood that an expanding, physical Universe could be realised without them.

Shortly thereafter, came Heckmann's popular realisation that expanding universes can be described with all values of $K$ and $\Lambda$. In particular, my previous statement that Heckmann's was the first treatment of the entire gamut of homogeneous, isotropic models within the expanding Universe paradigm is now justified; for though Robertson had found these solutions prior to Heckmann's paper, \cite{Heckmann1931}, he did not attempt to expose them as potential cosmological models. Heckmann had thus been the first to understand that both $K$ and $\Lambda$ might be zero in the expanding Universe, as one of nine general classifications of these geometries. Thus, his papers prompted Einstein and de~Sitter, in March 1932, to propose an idealised dust model for the Universe which completely neglects $\Lambda$, and would further neglect $K$ as a first approximation \cite{Einstein1932}. This is known as the Einstein-de~Sitter model, and was popularly thought to approximate the correct Universal geometry during the remainder of the twentieth century.

We should wonder why Einstein and de~Sitter did not give credit to Friedman, rather than Heckmann, for this suggestion, because the former had considered general $\Lambda$, and the case of zero curvature is a simple limit of his solutions. The important leap which they credited to Heckmann, was that he had pointed out `that the non-static solutions of the field equations of the general theory of relativity with constant density do not necessarily imply a positive curvature of three-dimensional space, but that this curvature may also be negative or zero.' But Friedman had found the negative curvature solution already in 1924, so this is obviously not true. In fact, neither Einstein nor de~Sitter had been aware of Friedman's second paper, \cite{Friedmann1924,Friedmann1999b}. In 1933 \cite{deSitter1933a}, de~Sitter wrote,
\begin{quotation}
Friedmann discusses the solutions of the field equations for different values of $\Lambda$. Lema\^{i}tre considers only positive $\Lambda$. Both authors consider a positive curvature of space only. The fact that both $\Lambda$ and the curvature may as well be negative or zero was only pointed out by Dr. Heckmann in July 1931.
\end{quotation}
The book from which this is taken, \textit{The Astronomical Aspect of the Theory of Relativity}, was published only one year before de~Sitter's death, so he likely never knew of Friedman's second paper; but Einstein may also never have known of it. In his writings on the cosmological problem following the paradigm shift in 1929 -- 1930---e.g., in \cite{Einstein1931}, as well as in his treatment of the problem in an appendix to \textit{The Meaning of Relativity} which he had published in 1945 \cite{Einstein1945}---Einstein only references Friedman's first paper \cite{Friedmann1922,Friedmann1999a} when crediting him with the original analytical discovery of the expanding Universe. 

The situation is yet more confusing when we note that Robertson had known of both papers by Friedman in 1929 \cite{Robertson1929}, when he had favoured de~Sitter's universe, and before the first one had been publicised by de~Sitter; and that in \cite{deSitter1933a}, in the paragraph \textit{directly preceding} the one which contains the above quotation, de~Sitter actually references \cite{Robertson1929} in connection with the proof that the Einstein and de~Sitter universes `are the only possible static, homogeneous, and isotropic solutions with positive curvature'. A further level of confusion comes in a review on the various descriptions of cosmology, written by Robertson in 1933 \cite{Robertson1933}, in which he notes that de~Sitter did write an article in 1932 \cite{deSitter1932}, where he performed an `analysis and classification of all of Friedman's worlds $p=0$.' However, in this article de~Sitter again only references Friedman's original paper \cite{Friedmann1922,Friedmann1999a}, citing Heckmann's paper \cite{Heckmann1931} as the first suggestion of the nine possible expanding worlds. Although Robertson's statement is technically correct, it is misleading because he had not noticed that de~Sitter had been unaware of the second paper, \cite{Friedmann1924,Friedmann1999b}.

The prevalent concept in all of the proposed descriptions of the Universe discussed in this section, is that of a separation of spacetime into homogeneous, isotropic spatial slices, comoving in proper time. This is due to the observation that the Universe appears to us to be isotropic, and to the cosmological principle---the idea that any observer in any galaxy must make essentially the same observations of the cosmos. The further assumption, that such a time exists, is what Bondi called `Weyl's postulate' \cite{Bondi1961}, although this is more restrictive than Weyl's original proposal, which had to do with a more general notion of causal coherence \cite{Ehlers2009}. Working from essentially these assumptions, Arthur Walker \cite{Walker1935a,Walker1935b,Walker1937} and Robertson \cite{Robertson1935,Robertson1936a,Robertson1936b} rigorously proved that the only kinematical systems which can exist can be represented by the three classes of isotropic, homogeneous spacetime geometries, 
\begin{equation}
ds^2=-dt^2+a(t)^2\left[\frac{dr^2}{1-Kr^2}+r^2(d\theta^2+\sin^2\theta{d}\varphi^2)\right],\label{RW}
\end{equation}
for which the factor $a(t)$ is conventionally scaled so that the curvature $K$ can be either $+1$, 0, or $-1$; i.e.~the three non-static solutions Robertson had found in \cite{Robertson1929}. These metrics are therefore called the Robertson-Walker (RW) metrics, and the cosmological models they describe are called Friedman-Lema\^{i}tre-Robertson-Walker (FLRW) models for the relevant contributions all four made to the development of the theory.

On the above assumptions, Robertson wrote, in 1929 \cite{Robertson1929},
\begin{quotation}
\ldots we have described the actual world in terms of co\"{o}rdinates which effect a natural separation of it into space and time and have determined its ideal background by a single assumption \ldots which is but the concrete expression of a uniformity implied by the very concept of a system of relativistic cosmology; in requiring that this ideal background can be fitted into the actual universe in the way described we have expressed in another way the assumption made by other writers on the subject, above all H. Weyl, that the world lines of all matter in the universe form a coherent pencil of geodesics.
\end{quotation}

Thus, Robertson had realised that this `natural separation' of space and time is not exactly what Weyl had intended. But regardless, this assumption, which may have been first explicitly stated by Robertson, and \textit{is} anyhow related to Weyl's postulate, has remained an elementary facet of standard relativistic cosmology, as a multitude of FLRW models have been used to describe popular cosmological theories ever since. As such, the general formulation of the cosmological problem has commonly been to define, at the outset, this division of spacetime into space and one Universal time, applicable both as the proper time of the Universe as a whole, and as that of every cosmic rest-frame observer, simultaneously; then, due to astronomical observations and the cosmological principle, to assume these spatial sections must be homogeneous and isotropic.

In \S~\ref{sec_MC}, a cosmological model is described which satisfies the cosmological principle, observed isotropy, and, in a more general sense, Weyl's prostulate, but which is geometrically dissimilar from the RW spacetimes, and thus disproves the common assertion that these geometries alone satisfy the principles of relativistic cosmology; and it serves to illustrate the error of Robertson's `natural separation' of spacetime, as, in that case, the cosmic standard time corresponding to Weyl's `assumption on the ``state of rest'' of stars' \cite{Weyl2009}, is clearly not a synchronous attribute of the proper times of each of those stars. However, the reasonable advantage of this model cannot be appreciated without significant theoretical development; so, we continue our description of the development of cosmological theory, in which we will rediscover the more natural and rational interpretation of cosmic expansion which was initially expounded by Eddington, which is formally inconsistent with FLRW big bang cosmology.

\section{Interpreting Cosmic Expansion\label{sec_ICE}}
Following Hubble's observational results, there emerged two schools of thought regarding $\Lambda$, and the one that won out can be generally associated with the notion, that `$\Lambda$ is \textit{not necessary} as long as the Universe is not static'. Although many words were used over the years, becoming more severe as time went on and the paradigm solidified, `unnecessary' seems to be the best to describe this interpretation, and was anyhow the original idea for which it became accepted. This interpretation is called the Einstein interpretation, for although it was first suggested by Friedman \cite{Friedmann1922,Friedmann1999a}, then Heckmann \cite{Heckmann1931}, and was favoured by many others, it had its strongest support from Einstein, who completely rejected the $\Lambda$-term and strongly promoted the $\Lambda\equiv0$ philosophy throughout the rest of his life. Many authors, e.g.~\cite{Pais1982,Chandrasekhar1983,Weinberg1989,Padmanabhan2003}, have written accounts of Einstein's interpretation of $\Lambda$, but it is worthwhile to produce a full exposition of this idea here because many of the others are either brief, misleading, or even incorrect.

Scepticism towards Einstein's introduction of the $\Lambda$-term in 1917 \cite{Einstein1917} came almost immediately. When de~Sitter proposed his solution later that year \cite{deSitter1917}, he wrote, 
\begin{quotation}
It cannot be denied that the introduction of the constant $\Lambda$\ldots is somewhat artificial, and detracts from the simplicity and elegance of the original theory of 1915, one of whose great charms was that it embraced so much without introducing any new empirical constant.
\end{quotation}
As an astronomer, de~Sitter had however been confronted with a necessity to account for the redshifts of spiral nebulae throughout the 1920s; thus, his own sentiments towards the $\Lambda$-term had become less severe by the time of Hubble's discovery, due to Eddington's interpretation of a cosmic dispersal in a de~Sitter cosmology. In the majority of de~Sitter's articles from 1930 until his death in 1934 (see, e.g., \cite{deSitter1930,deSitter1933a,deSitter1933b}), with the exception of the one he co-authored with Einstein in 1932 \cite{Einstein1932}, he had invoked the $\Lambda$-term to describe the present nature of the expansion inferred from Hubble's results. But although de~Sitter's initial reaction towards $\Lambda$ had become less severe due to Eddington's and Hubble's results, the influence of such sentiments as the one he initially expressed would be everlasting on Einstein, who, in his later writings on the cosmological problem, would emulate de~Sitter's initial misgivings.
 
Einstein's distaste for the $\Lambda$-term is already evidenced in a postcard he sent to Weyl on May 23, 1923 [Einstein Archives call number [24-81.00]; original held by the Weyl-Archiv at the Eidgen\"{o}ssische Technische Hochschule Z\"{u}rich], following the publication of Weyl's cosmological theory. In this letter, he wrote,
\begin{quote}
[Des kosmologischen Gliedes] w\"{u}rde zun\"{a}chst nichts schaden, weil man die Theorie ganz zwanglos verallgemeinern kann, was ich unterdessen gefunden habe.\ldots Ich sende Ihnen dann die Korrektur meiner Arbeit, die ich unbedingt publizieren muss, weil der Eddington'sche Gedanke notwendig zu Ende gedacht werden muss. Ich glaube jetzt auch, dass alle diese Versuche auf rein formaler Basis die physikalische Erkenntnis nicht weiter bringen werden. Vielleicht hat die Feldtheorie schon alles hergegeben, was in ihren M\"{o}glichkeiten liegt. Inbezug auf das kosmologische Problem bin ich nicht Ihrer Meinung. Nach de~Sitter laufen zwei gen\"{u}gend von einander entfernte materielle Punkte beschleunigt auseinander. Wenn schon keine quasi-statistische Welt, dann fort mit dem kosmologischen Glied. [[The cosmological term] would not hurt at first, because one can generalise the theory quite naturally, as I have meanwhile found.\ldots I send you then the correction of my work, which I absolutely must publish because the Eddington-ish considerations must necessarily be thought through. I now also believe that all these attempts towards a purely formal basis will not lead to further physical knowledge. Perhaps the field theory has already given away everything that lies in its possibilities. With reference to the cosmological problem, I am not of your opinion. Following de~Sitter, we know that two sufficiently separate material points are accelerated from one another. If there is no quasi-static world, then away with the cosmological term.]
\end{quote}
Eight days later, Einstein submitted his second note on Friedman's paper \cite{Einstein1923}, which he now referred to as `correct and clarifying'. His use of the word clarifying (`aufkl\"{a}rend') in the context of a retraction of his previous statement should have seemed puzzling without this letter to Weyl; however, now that we have it in the context of the other interpretations of expansion in de~Sitter spacetime by Eddington and Weyl, it is clear that Einstein had been expressing his preference for Friedman's solution as a possible way out of the alternate description of accelerating expansion. 

But Friedman's solution had not been known to everyone, and Weyl's confusion over this letter is quite plain in a dialogue he wrote the following year, on inertial mass and the cosmos, in which an orthodox `Petrus' (St.~Peter), who is supposed to represent Einstein's position, eventually responds to a more radical---in fact, as Weyl must have seen things, dangerously close to apostatical\footnote{At a point early on in the dialogue, `Paulus' actually says, `wenn hier der Fels liegt, auf dem die Relativit\"{a}tskirche steht, o Petrus!, so bin ich in der Tat ein Abtr\"{u}nniger geworden. Aber um dich \"{u}ber meine Ketzerei ein wenig zu beruhigen, \ldots' (if here lies the rock on which the relativity church stands [cf.~Matthew~16:18], o Peter!, then I have indeed become an apostate. But to calm you a little about my heresy, \ldots) \cite{Weyl1924}.}---`Paulus' (St.~Paul), who presents Weyl's ideas \cite{Weyl1924},
\begin{quote}
Wenn es mit Hilfe des kosmologischen Gliedes nicht gelingt, das Machsche Prinzip durchzuf\"{u}hren, so halte ich es \"{u}berhaupt f\"{u}r zwecklos und bin f\"{u}r die R\"{u}ckkehr zur elementaren Kosmologie. [If the cosmological term fails to help with leading through to Mach’s principle, then I consider it to be generally useless, and am for the return to the elementary cosmology.]
\end{quote}
`Petrus' was speaking here of a return to a special relativity-based cosmology; and `Paulus' went on to reason that this would be too hasty, explaining that the expanding `de~Sitter cosmology' would account better for the observational data than either this `elementary cosmology' or Einstein's statical model. Weyl obviously was, like nearly everyone else in 1924, completely unaware of Friedman's original solution, and it would seem oddly devious if Einstein, who had apparently been pleased with the clarification that had been afforded by Friedman's solution, actually made no subsequent attempt to clear up this matter by sharing what he knew. 

But regardless of the way circumstances played out, it is clear that Einstein had been unsatisfied with the $\Lambda$-term throughout most of the 1920s, and it is therefore not surprising that when he submitted his paper on the cosmological problem in 1931 \cite{Einstein1931} he would attempt to retract it. He mentions the $\Lambda$-term with negative connotations twice in this article, from which we see hints of de~Sitter's speculation. After giving an introduction to the new picture of the expanding Universe, Einstein writes,
\begin{quotation}
Unter diesen Umst\"{a}nden [\textsc{Friedmanns} und \textsc{Hubbels} Resultate] mu{\ss} man sich die Frage vorlegen, ob man den Tatsachen ohne die einf\"{u}hrung des theoretisch ohnedies unbefriedigenden $\Lambda$-Gliedes gerecht werden kann. [With these understandings [Friedman's and Hubble's results], one must ask oneself whether the facts can be accounted for without the introduction of the (theoretically, anyhow) unsatisfactory $\Lambda$-term.]
\end{quotation}
Then, in conclusion,
\begin{quote}
\ldots Bemerkenswert ist vor allem, da{\ss} die allgemeine Relativit\"{a}tstheorie \textsc{Hubbels} neuen Tatsachen ungezwungener (n\"{a}mlich ohne $\Lambda$-Glied) gerecht werden zu k\"{o}nnenscheint als dem nun empirisch in die Ferne ger\"{u}ckten Postulat von der quasistatischen Natur des Raumes. [Above all, it is remarkable that Hubble's new facts allow general relativity theory to seem less contrived (namely, without the $\Lambda$-term), as the postulate of the quasistatic nature of space has moved into the distance.] 
\end{quote} 
Thus, Einstein had thought of the $\Lambda$-term as a blemish to his theory, which had been added \textit{ad hoc}, in order to account for an incorrect world-view. With these remarks, he seems glad to be rid of the term which had not only itself seemed contrived, but which he feared had made general relativity theory seem contrived, allowing its critics an argument which could now be revoked.

A year later, Einstein and de~Sitter published their cosmological model \cite{Einstein1932}, in which $\Lambda$ was to be neglected from the outset. The treatment of $\Lambda$ in this article was mis-interpreted by Subrahmanyan Chandrasekhar \cite{Chandrasekhar1983}, who wrote,
\begin{quote}
\ldots When Friedmann's cosmological models were found to provide an adequate base for accounting for the simple fact of the Hubble expansion, Einstein and de~Sitter, in a joint paper, stated that one can do without the cosmical constant. In view of the many exaggerated statements that have been made concerning this supposed `retraction' of $\Lambda$, it is of interest to record precisely what it was they said.
\begin{quotation}
Historically the term containing the ``cosmological constant'' $\Lambda$ was introduced into the field equations in order to enable us to account theoretically for the existence of a finite mean density in a static universe. It now appears that in the dynamical case this end can be reached without the introduction of $\Lambda$\ldots

\ldots The curvature [constant $\Lambda$] is, however, essentially determinable, and an increase in the precision of data derived from observations will enable us in the future to fix its sign and to determine its value.
\end{quotation}
\end{quote}
The edit in the second paragraph was added by Chandrasekhar in order to support his argument that Einstein and de~Sitter had not rejected $\Lambda$. However, it is absolutely clear in the two-page article \cite{Einstein1932}, that the curvature of which they speak is not $\Lambda$, but the curvature of three-dimensional space, $K$; and it is shocking that Chandrasekhar did not understand this. Regardless, the cosmological constant was neglected altogether in the Einstein-de~Sitter model because in the full FLRW cosmological models $\Lambda$ truly is unnecessary in order to account for the observed expansion, as Friedman and Heckmann had previously noted.

So, as of 1932, Einstein had (jointly) proposed a model of the Universe which rejected the $\Lambda$-term, which he had previously thought to be theoretically unsatisfactory. These thoughts appear to have strengthened in subsequent years, and in 1945 he added an appendix to \textit{The Meaning of Relativity} \cite{Einstein1945}, in which he includes many damning remarks towards the existence of $\Lambda$; for example,
\begin{quotation}
\ldots The introduction of [$\Lambda$] constitutes a complication of [general relativity] theory, which seriously reduces its logical simplicity. Its introduction can only be justified by the difficulty produced by the almost unavoidable introduction of a finite average density of matter. We may remark, by the way, that in Newton's theory there exists the same difficulty.

The mathematician Friedman found a way out of this dilemma.\symbolfootnote[1]{He showed that it is possible, according to the field equations, to have a finite density in the whole (three-dimensional) space, without enlarging the field equations \textit{ad hoc}. Zeitschr.~f.~Phys.~10 (1922). [Einstein's footnote.]} His result then found a surprising confirmation by Hubble's discovery of the expansion of the stellar system.
\end{quotation}
And then, in a summary note, 
\begin{quotation}
The introduction of the ``cosmologic member'' into the equations of gravity, though possible from the point of view of relativity, is to be rejected from the point of view of logical economy. As Friedman was the first to show one can reconcile an everywhere finite density of matter with the original form of the equations of gravity if one admits the time variability of the metric distance of two mass points.\symbolfootnote[1]{If Hubble's expansion had been discovered at the time of the creation of the general theory of relativity, the cosmologic member would never have been introduced. It seems now so much less justified to introduce such a member into the field equations, since its introduction loses its sole original justification,---that of leading to a natural solution of the cosmologic problem. [Einstein's footnote]}
\end{quotation}

Eventually, Einstein's overall interpretation of $\Lambda$ was appropriately summarised by Wolfgang Pauli in his 1958 supplemental notes to the English addition of \textit{Theory of Relativity} \cite{Pauli1958}, which had originally been published in German in 1921:
\begin{quotation}
Einstein was soon aware of [the results of Friedman, Lema\^{i}tre, and Hubble] and \textit{completely rejected the cosmological term} as superfluous and no longer justified. I fully accept this new standpoint of Einstein's.
\end{quotation}

Finally, in 1970, when George Gamow's autobiography was published posthumously, we learned that `much later [than Hubble's discovery], when [Gamow] was discussing cosmological problems with Einstein, he remarked that the introduction of the cosmological term was the biggest blunder he ever made in his life' \cite{Gamow1970}. Now, whether this was considered a blunder because he may otherwise have discovered Friedman's solution and predicted Hubble's result; or because now that he had suggested the term it would never go away completely, cropping up, e.g., as an integral component in Eddington's and Lema\^{i}tre's models, and later in the steady state theory; or because of something else entirely, we'll never know. In fact, we can never truly know whether these \textit{were} his words because they came second hand, as well as posthumously on both accounts! 

However the true statement may have been made or meant, this completes our comprehensive exposition of Einstein's interpretation of $\Lambda$ in light of the observed cosmic expansion. In his own words, Einstein did abandon the cosmological term unequivocally, and Pauli's summary of Einstein's view seems appropriate. Although the description of this view was given gradually, over many years and with increasingly severe tones, it always remained mathematically equivalent to the Einstein-de~Sitter model, proposed in 1932 \cite{Einstein1932}. And as may have been expected when the master of the theory had cast his lot, the \textit{status quo} went with him. Thus, the Einstein-de~Sitter model, effectively parametrised by dust-density and curvature, became the standard model in cosmology throughout most of the twentieth century---and it was for many a great shock to find that the Universe is otherwise.

This should not have been so; for there is a natural explanation of expansion which was popularly neglected, although many of the prominent researchers in early cosmology---most notably, Eddington---had tried to make it work. This interpretation, which I have briefly mentioned above, is that the cosmic expansion is driven---and in fact has been for all time---by positive $\Lambda$. There are difficulties with this `Eddington interpretation' which have always been known; most notably, that, with the exception of Weyl's model, it cannot be accounted for in any FLRW big bang universe. For as we follow Friedman's equations, Eqs.~(\ref{Friedmann}}) -- (\ref{Friedmann2}}), backwards in time, the constant $\Lambda$ must inevitably become negligible when compared with the increasing density of matter. Accordingly, the Einstein interpretation is the only interpretation of cosmic expansion that has existed in the literature since the discovery of the cosmic microwave background (CMB) in 1965, which provided strong confirmation of the Big Bang theory.

Therefore, the overwhelming evidence for positive $\Lambda$, demonstrated from observations of type Ia supernovae (SNe Ia) by two separate groups at the end of the last century \cite{Riess1998,Perlmutter1999}, allowed for only a superficial contrast between the interpretations of cosmic expansion and observational results. Furthermore, quantum theory had already predicted a non-zero vacuum energy which could act like the observed cosmological constant at late times, after having evolved from a much larger value early on. Since the discovery of non-zero $\Lambda$, there has been much hope for statistical evidence of such a \textit{quintessence} field, which would decay in such a manner, and could thus be attributed as a solution to the \textit{cosmological constant problem}---the discrepancy between values of the vacuum energy predicted by quantum field theory and observed $\Lambda$---and a suggestive link to inflation theory; see, e.g.,~\cite{Weinberg2008} for a recent comprehensive treatment of this picture. This prediction, along with the fact that $\Lambda$ has no essential role in the standard cosmology, resulted, almost immediately after its discovery, in re-naming the observed cosmological constant, rather as `dark energy',---a blanket term to cover all theories.

In the remainder of this chapter, I will continue the historical account of cosmology and the cosmological constant from the Eddington interpretation onwards, and eventually argue against the above theories in that context. The result of this exposition will be that the standard cosmology, which has been increasingly refined since the 1930s, and is now known to be wrought with many problems---e.g.~the flatness, horizon, and cosmological constant problems which the quintessence and inflation theories aim to resolve---should likely need to be abandoned in favour of one that would account \textit{a priori} for the facts of observation. 

To begin, let us ignore the evidence of the Big Bang---which will anyhow be shown to have a description in the new cosmology of Chapter~\ref{chap_ANC}---so that the full argument may be presented in an historical context, without having to accept that the new cosmological model, which is consistent with this interpretation, does work. Then we are perfectly at liberty, as Eddington was, to argue the possibility of his interpretation of $\Lambda$, which contends that \textit{positive $\Lambda$ must be the sole cause of cosmic expansion}. 

Knowing only that the Universe is expanding at a rate which, near to us, looks like a linear velocity-distance relation, the most rational interpretation certainly is that the expansion is driven by positive $\Lambda$ at all times, rather than as the result of a singular beginning. For the latter can only be described with a full understanding of the process of that beginning, which itself must have been very special in order to account for the uniform expansion we observe everywhere (see, e.g., the discussion of this point in \cite{Penrose2005,Penrose2009}), and anyhow cannot be explained with the FLRW model alone, while the former is a known facet of general relativity theory, which Einstein had to admit was `possible from the point of view of relativity', and which has the possibility of naturally explaining the observed expansion. 

Eddington's belief in the $\Lambda$-term did not have the luxury of observational confirmation, as ours now does. He believed $\Lambda$ was a fundamental constant of the Universe on `general philosophical grounds', which he had held before and apart from the observed expansion, as mentioned in \S~\ref{sec_EMC} \cite{Eddington1923,Eddington1929,Eddington1931}. But whether or not Eddington's reasons for believing in $\Lambda$ as a fundamental facet of general relativity theory are ultimately correct, the simple fact remains that $\Lambda$ is \textit{not} superfluous or unjustified in the mathematical theory, and should not be rejected at the outset. The case $\Lambda=0$ is very special indeed, and setting $\Lambda\equiv0$ `from the point of view of logical economy', so as not to complicate the logical simplicity of the theory, is wrong, and constitutes a breach of the most natural extension of the cosmological principle; i.e.~that our \textit{Universe} should not be very special. From this perspective alone, the most logical interpretation of the observation that the Universe is expanding, is that our Universe possesses a positive cosmological constant which drives the expansion; for this is the only interpretation which would go beyond the standard empirical model and provide a natural theoretical \textit{explanation} of it.

In 1933, Eddington published \textit{The Expanding Universe} \cite{Eddington1933}, in which he fully described his interpretation of the cosmic expansion. Although his opinions are often extreme, and will require a little further comment, his idea is made most clear through a string of direct quotations from this book, which concisely represents the picture he held of the Universal expansion:
\begin{quotation}
The \textit{spiral nebulae} \ldots are extra-galactic objects; that is to say, they lie beyond the limits of the Milky Way aggregation of stars which is the system to which our sun belongs, and are separated from it by wide gulfs of empty space. \ldots

Each island system is believed to be an aggregation of thousands of millions of stars with a general resemblance to our own Milky Way system. As in our own system there may be along with the stars great tracts of nebulosity, sometimes luminous, sometimes dark and obscuring. \ldots

The lesson of humility has so often been brought home to us in astronomy that we almost automatically adopt the view that our own galaxy is not specially distinguished---not more important in the scheme of nature than the millions of other island galaxies. \ldots

When the collected data as to radial velocities and distances [of these galaxies] are examined a very interesting feature is revealed. The velocities are large, generally very much larger than ordinary stellar velocities. The more distant nebulae have the bigger velocities; the results seem to agree very well with a linear law of increase, the velocity being simply proportional to the distance. The most striking feature is that the galaxies are almost unanimously running away from us. \ldots

We can exclude the spiral nebulae which are more or less hesitating as to whether they shall leave us by drawing a sphere of rather more than a million light-years radius round our galaxy. \textit{In the region beyond, more than 80 have been observed to be moving outwards, and not one has been found coming in to take their place}.

The inference is that in the course of time all the spiral nebulae will withdraw to a greater distance, evacuating the part of space that we now survey. Ultimately they will be out of reach of our telescopes unless telescopic power is increased to correspond.\footnote{This speculation has been disproven by Lawrence Krauss and Glenn Starkman \cite{Krauss2000}. As it turns out, there is no chance of observing light from distant galaxies after a finite time in our future, as the Universe asymptotically approaches de~Sitter spacetime.} \ldots

The unanimity with which the galaxies are running away looks almost as though they had a pointed aversion to us. We wonder why we should be shunned as though our system were a plague spot in the universe. But that is too hasty an inference, and there is really no reason to think that the animus is especially directed against our galaxy. If this lecture room were to expand to twice its present size, the seats all separating from each other in proportion, you would notice that everyone had moved away from you. Your neighbour who was 2 feet away is now 4 feet away; the man over yonder who was 40 feet away is now 80 feet away. It is not \textit{you} they are avoiding; everyone is having the same experience. \ldots

\ldots the picture of the universe now in the minds of those who have been engaged in practical exploration of its large-scale features \ldots is the picture of an \textit{expanding universe}. The super-system of the galaxies is dispersing as a puff of smoke disperses. Sometimes I wonder whether there may not be a greater scale of existence of things, in which it \textit{is} no more than a puff of smoke. \ldots

\ldots I want to speak of the alteration that Einstein made in his law of gravitation. The amended law is written $G_{\mu\nu}=\Lambda{g_{\mu\nu}}$, and contains a natural constant $\Lambda$ called the \textit{cosmical constant}. The term $\Lambda{g_{\mu\nu}}$ is called the \textit{comsical term}. The constant is so small that in ordinary applications to the solar system, etc., we set it equal to zero, and so revert to the original law $G_{\mu\nu}=0$. But however small $\Lambda$ may be, the amended law presents the phenomenon of gravitation to us in a new light, and has greatly helped to an understanding of its real significance; moreover, we have now reason to think that $\Lambda$ is not so small as to be entirely beyond observation. The nature of the alteration can be stated as follows: the original law stated that a certain geometrical characteristic ($G_{\mu\nu}$) of empty space is always zero; the revised law states that it is always in a constant ratio to another geometrical characteristic ($g_{\mu\nu}$). We may say that the first form of the law utterly dissociates the two characteristics by making one of them zero and therefore independent of the other; the second form intimately connects them. Geometers can invent spaces which have not either of these properties; but actual space, surveyed by physical measurement, is not of so unlimited a nature.

We have already said that the original term in the law gives rise to what is practically the Newtonian attraction between material objects. It is found similarly that the added term ($\Lambda{g_{\mu\nu}}$) gives rise to a repulsion directly proportional to the distance. Distance from what? Distance from \textit{anywhere}; in particular, distance from the observer. It is a dispersive force like that which I imagined as scattering apart the audience in the lecture-room. Each thinks it is directed away from him. We may say that the repulsion has no centre, or that every point is a centre of repulsion.

Thus in straightening out his law of gravitation to satisfy certain ideal conditions, Einstein almost inadvertently added a repulsive scattering force to the Newtonian attraction of bodies. We call this force the \textit{cosmical repulsion}, for it depends on and is proportional to the cosmical constant. It is utterly imperceptible within the solar system or between the sun and neighbouring stars. But since it increases proportionately to the distance we have only to go far enough to find it appreciable, then strong, and ultimately overwhelming. In practical observation the farthest we have yet gone is 150 million light-years. Well within that distance we find that celestial objects are scattering apart as if under a dispersive force. Provisionally we conclude that here cosmical repulsion has become dominant and is responsible for the dispersion. \ldots

\ldots [The cosmical constant] renders the theory of gravitation and its relation to space-time measurement so much more illuminating, and indeed self-evident, that to return to the earlier view is unthinkable. I would as soon think of reverting to Newtonian theory as of dropping the cosmical constant. \ldots

\ldots There are only two ways of accounting for large receding velocities of the nebulae: (1) They have been produced by an outward directed force as we here suppose, or (2) as large or larger velocities have existed from the beginning of the present order of things.\symbolfootnote[1]{For completeness we must add the possible hypothesis that the system once extended much further than now, that it collapsed, and is now on the rebound. This allows the large velocities to have been produced by \textit{inward} directed force, the inward velocities being turned into outward velocities by passage through the centre. So far as I know, this is not advocated by anyone. It does not seem capable of providing for the distribution of velocities which we observe. [Eddington's footnote.]} Several rival explanations of the recession of the nebulae, which do not accept it as evidence of the repulsive force, have been put forward. These necessarily adopt the second alternative, and postulate that the large velocities have existed from the beginning. This might be true; but can scarcely be called an \textit{explanation} of the large velocities. \ldots

\ldots the theory recently suggested by Einstein and de~Sitter, that in the beginning all the matter created was projected with a radial motion so as to disperse even faster than the present rate of dispersal of galaxies,\symbolfootnote[1]{They do not state this in words, but it is the meaning of their mathematical formulae. [Eddington's footnote.]} leaves me cold. One cannot deny the possibility, but it is difficult to see what mental satisfaction such a theory is supposed to afford. \ldots

When once it is admitted that there exists everywhere a radius of curvature ready to serve as comparison standard, and that spatial distances are directly or indirectly expressed in terms of this standard, the law of gravitation ($G_{\mu\nu}=\Lambda{g_{\mu\nu}}$) follows without further assumption; and accordingly the existence of the cosmical constant $\Lambda$ with the corresponding force of cosmical repulsion is established. Being in this way based on a fundamental necessity of physical space, the position of the cosmical constant seems to me impregnable; and if ever the theory of relativity falls into disrepute the cosmical constant will be the last stronghold to collapse. \textit{To drop the cosmical constant would knock the bottom out of space}. \ldots
\end{quotation}

Now, whether $\Lambda$ should be so \textit{required} in the physical description of gravity by general relativity theory as Eddington argues, is not yet certain. Indeed, Chandrasekhar \cite{Chandrasekhar1983} commented that `no serious student of relativity is likely to subscribe to Eddington's view that ``to set $\Lambda=0$ is to knock the bottom out of space''.' And Eddington's statement, that he would `as soon think of reverting to Newtonian theory as of dropping the cosmical constant' is obviously exaggerated; for he was intimately aware of the important confirmations the theory had received in describing the perihelion precession of Mercury and the deflection of light,---though he may have been as likely to revert to the theory of modified Newtonian dynamics (MOND) proposed by Mordehai Milgrom \cite{Milgrom1983a,Milgrom1983b,Milgrom1983c} as general relativity theory with $\Lambda\equiv0$, had he been alive in 1983. To Eddington, the theory had only really been complete with the introduction of positive $\Lambda$, which, he thought, serves as a physical basis for the length scales the theory would describe. But whether or not anyone would currently subscribe to Eddington's view, to simply dismiss the notions of one of the great natural philosophers of the twentieth century concerning a concept which he had given such considerable thought, without giving the problem an equivalent amount of thought, seems too hasty an act. Thus, I note only that on purely aesthetic or philosophical grounds we can currently make no headway on the matter: for two of the primary exponents of the theory, Einstein and Eddington, had extreme contrasting views. However, we don't need this in order to proceed with Eddington's interpretation of \textit{cosmic expansion}, so it can be left for now as an open issue. All that is required is that $\Lambda\neq0$ should necessarily be considered a possibility of the mathematical theory of relativity. On this point, no serious student of relativity would currently object, as from a mathematical standpoint Pauli's excessive statements regarding $\Lambda$ are incorrect!

Working from only this more moderate view and the observation that the Universe is uniformly expanding, Eddington's theory \cite{Eddington1930}, which was anticipated by Lema\^{i}tre \cite{Lemaitre1927,Lemaitre1931a}, in which the Einstein Universe would expand \textit{due to} $\Lambda$ was far more natural than the theory that would come to dominate. Eddington's interpretation \textit{explains} the Universal expansion \textit{at all times}, whereas the Einstein interpretation \textit{never} does---for, even if we would set additional conditions in the very early Universe that would contribute to the subsequent physical state of expansion, the inexplicable singularity still remains the initial cause, according to the mathematical solution. However, in order to allow for the implicit expansion of space that is so naturally accounted for by positive $\Lambda$, Eddington had to relinquish any possibility of a beginning of expansion, which he regarded anyhow (or perhaps because it did not fit with his interpretation) as `too unaesthetically abrupt'. As mentioned above, this is true for any FLRW big bang universe aside from Weyl's `de~Sitter cosmology': although expansion may be driven by $\Lambda$ at late times, it cannot account for the expansion shortly after the big bang, which must be caused by the big bang singularity itself.

\section{Towards a Standard Cosmogony\label{sec_TSC}}
In light of the problem Eddington's interpretation has in reconciling itself with the FLRW big bang models, many of the other theorists working on the problem, e.g.~Lema\^{i}tre and de~Sitter, would take the more moderate stance, that $\Lambda\neq0$ is perfectly allowable by general relativity theory, and could possibly have a fundamental role even, so that, despite any uncertainty in that regard, it should always be included in general cosmological descriptions. I've mentioned previously, that, following Hubble's discovery, de~Sitter's only apparent deviation from general descriptions with $\Lambda\neq0$ appears in the paper he wrote with Einstein in 1932 \cite{Einstein1932}---otherwise, his articles always accounted for this generality; although, in conclusion to the comprehensive report he wrote, containing his first reaction after learning of Lema\^{i}tre's article \cite{Lemaitre1927} from Eddington, he had clearly not come to the whole-hearted acceptance that Eddington had \cite{deSitter1930}:
\begin{quote}
The constant $\Lambda$, which is a measure of the inherent expanding force of the universe, is still very mysterious, and it is difficult to see what its real meaning is. It might even be thought to be one constant too many, unless we may hope that it will ultimately be found to be in some way connected with Planck's constant $h$.
\end{quote}
And, in \cite{Chandrasekhar1983}, Chandrasekhar wrote that he 
\begin{quote}
once asked Lema\^{i}tre, sometime during the late fifties, what, in his judgement, was the most important change wrought by the general theory of relativity in our basic physical concepts. His answer, without a moment's hesitation, was `the introduction of the cosmical constant!'
\end{quote}
Indeed, Lema\^{i}tre appears always to have been influenced by his time spent at Cambridge, and in 1949 he described `the cosmical constant as a second constant demanded by the logical structure of [general relativity] theory' \cite{Lemaitre1949}. 

Lema\^{i}tre's adherence to the cosmological constant may in the end have come from an astronomical standpoint, as he believed throughout the remainder of his career that positive $\Lambda$ could account for the observed clustering of galaxies in the Universe \cite{Lemaitre1933,Lemaitre1997,Lemaitre1934,Lemaitre1949,Lemaitre1952,Lemaitre1958}, which remains a little-known---but likely somewhat valid---interpretation. Lema\^{i}tre first described his theory in 1933 \cite{Lemaitre1933}, in an analysis which determined the gravitational boundary of a massive cluster. This theory, which he later credited to de~Sitter, in \cite{Lemaitre1949}, who had dedicated a section of \cite{deSitter1933a} to an astronomical analysis of the `Balance of gravitation and expansive force', seems to have been commonly conceived due to the observed blueshifts of nearby galaxies; e.g., as Eddington described in the above quotation. 

Essentially, it was well-known that some nearby galaxies have blueshifted spectra, and therefore from our perspective they do not appear to take part in the greater cosmic expansion, but are gravitationally bound to the Local Group. Because observational evidence exists that such structures dominate the Universe---i.e., many galaxies appear to be members of clusters containing various numbers of galaxies---de~Sitter and Lema\^{i}tre proposed that a similar situation exists throughout the Universe, as in our Local Group, in which there is a boundary at the edge of each cluster where the `Newtonian' attraction is balanced by the cosmic repulsion due to $\Lambda$. In an idealised situation, in which the mass of the cluster is spherically symmetric and non-rotating, the radius of this boundary is simple to calculate, and Lema\^{i}tre used what is now referred to as the Lema\^{i}tre-Tolman \cite{Tolman1934} metric and found that the equilibrium radius of a spherical mass is
\begin{equation}
{r_e}^3=\frac{3GM}{\Lambda{c^2}},\label{r_e}
\end{equation}
or, adopting $\Lambda=1.27\times10^{-52}~\mathrm{m}^{-2}$, from \cite{Hinshaw2009}, along with modern values of the other constants,
\begin{eqnarray}
M\!\!\!&=&\!\!\!\nonumber{5.70\times10^{-26}~\mathrm{kg~m}^{-3}}\cdot{{r_e}^3},\\
\!\!\!&=&\!\!\!\left(\frac{r_e}{106~\mathrm{pc}}\right)^3~\mathrm{M_{\odot}},\label{Lemaitre_mass}
\end{eqnarray}
where $1~\mathrm{M_{\odot}}$ is the Sun's mass.\footnote{Lema\^{i}tre adopted $\Lambda=10^{-50}~\mathrm{m}^{-2}$, and thus calculated a denominator of $(80~\mathrm{ly})^3$, or $(25~\mathrm{pc})^3$ \cite{Lemaitre1933}.} 

In 1954, Felix Pirani proposed a similar theory \cite{Pirani1954}: he found, equivalently, that the radius $r_e$ in Eq.~(\ref{r_e}) is that of the largest circular orbit about a SdS mass which a test-particle can maintain, and proposed that outside the corresponding sphere---i.e., Eddington's `sphere of rather more than a million light-years radius round our galaxy'---an orbiting object in any galaxy cluster would necessarily spiral outwards due to the cosmic repulsion. More detailed investigations of the SdS geometry, which include this result, are given without reference to astrophysical implications in \cite{Stuchlik1983,Stuchlik1999}.

The Local Group's equilibrium radius has recently been determined to be $r_{e_{LG}}=(0.96\pm0.03)~\mathrm{Mpc}$ \cite{Karachentsev2009}. Using this value in Lema\^{i}tre's mass equation, Eq.~(\ref{Lemaitre_mass}), we find $M_{LG}=(7.4\pm0.7)~10^{11}~\mathrm{M_{\odot}}$. The mass determined from this value of $r_{e_{LG}}$ in \cite{Karachentsev2009}, from an analysis which attempts to describe the boundary within the Einstein interpretation of cosmic expansion, according to which the Hubble flow would tear galaxies away from the perimeters of massive clusters---i.e.~according to an imbedding in the overall FLRW Universe, with its best determined parameters---is $M_{LG}=(1.9\pm0.2)~10^{12}~\mathrm{M_{\odot}}$.  Due to the vastly different methods from which these two results were obtained, and to the non-trivial assumptions involved in each analysis, this factor of $\lesssim3$ discrepancy is surprisingly small.

When Lema\^{i}tre further developed this theory, he noted a duality admitted by Einstein's equation: that $\Lambda$ can equivalently be thought of as a vacuum energy with negative pressure, with equation-of-state, $p_{\Lambda}=-\rho_{\Lambda}$ \cite{Lemaitre1934}. This interpretation---which was essentially anticipated by de~Sitter \cite{deSitter1917}, who wrote,
\begin{quote} 
In Einstein's theory of general relativity there is no difference between gravitation and inertia. The combined effect of the two is described by the fundamental tensor $g_{\mu\nu}$, and how much of it is to be called inertia and how much gravitation is entirely arbitrary. We might abolish one of the two words, and call the whole by one name only.
\end{quote} 
---helped to describe the theory Lema\^{i}tre would always adhere to, in which systems initially in equilibrium might be perturbed in such a manner that some regions would collapse, becoming separated ever further from each other by intermediate expanding space, driven by an effective negative pressure. A further development, along similar lines of thought, will be discussed in \S~\ref{sec_DPM}.

Historically, the most noted significance of Lema\^{i}tre's interpretation has been that it provided a link in general relativity theory---especially in the cosmological problem---to a developing concept from quantum theory:---that the vacuum may have non-zero energy; see, e.g., \cite{Rugh2002,Padmanabhan2003} for descriptions of the early history. In the simplest scenario of this theory, in which the vacuum energy would not evolve in time according to some quintessence field, its density, being independent of spatial expansion or contraction, would remain constant; so the vacuum would have an equation-of-state equivalent to that of $\Lambda$. Thus, Lema\^{i}tre had provided the analytical prediction, that, if the vacuum energy were non-zero, it might likely act as a cosmological constant.

At the very heart of Lema\^{i}tre's interpretation, we have therefore found an answer to de~Sitter's rhetoric: although gravity and inertia are formally equivalent in general relativity theory, there remains a basic difference between equivalence and equality, between relative perception and reality, which must always be considered if we care to understand the essential truth of any system described by the mathematical equations.

After these attempts to explain phenomena by positive $\Lambda$, few others were made in the remainder of the twentieth century. As well, the concept of vacuum energy would take an equally long time to re-emerge as a central problem of modern physics. In astronomy, $\Lambda$ would thus slowly fade into obscurity. 

The most notable resurrection of the cosmological constant prior to its observation was in the Steady State theory, proposed in 1948 by Bondi and Thomas Gold \cite{Bondi1948} and by Hoyle \cite{Hoyle1948}, which made use of $\Lambda$ in the same vein as in the Eddington interpretation, and was in fact closely related to Weyl's theory. The Steady State theory, like Weyl's cosmology, would describe the Universe as a slice of de~Sitter spacetime---i.e.~according to Weyl's postulate---with infinite age, in which the cosmic expansion would be determined solely by $\Lambda$. However, rather than having a defined beginning at $t=-\infty$ when all physical distances converge to zero, and exponentially expanding until our Galaxy would appear alone in the Universe after the signals from all others would be redshifted beyond detectability, as Weyl's model required, the Steady State theory operated from an ammended, `Perfect Cosmological Principle', in which the Universe would appear as it does to all observers in the Universe \textit{for all time}. In order to account for the observed galaxy distribution and expansion, the theory required a steady creation of matter from the vacuum, which would collapse to form galaxies in intergalactic space. It would thus avoid the aesthetical issue that in the infinite existence of the Universe there would be only a short time during which life could exist and the Universe would appear as it does, with the multitude of galaxies we observe, which is a noted downside to most other cosmological theories, especially those of Weyl and Eddington. 

This theory was however refuted when the CMB was discovered in 1965 by Arno Penzias and Robert Wilson \cite{Dicke1965,Penzias1965}, which confirmed a prediction of the theory on the origin of elemental abundances developed by Gamow, Ralph Alpher, and Robert Herman in the late 1940s \cite{Gamow1946,Alpher1948a,Gamow1948,Alpher1948b,Alpher1950a}. This theory would account for the abundances of the light elements in the Universe---most notably, the roughly 25\% abundance of helium---as well as predict the existence of a cosmic background radiation (CBR) produced in the initial hot state when the Universe had sufficiently cooled so that radiation could decouple from matter, allowing electrons to recombine (for the first time) with nuclei, which would subsequently have cooled through the expansion of space to `about $5^{\circ}$~K' \cite{Alpher1948b}. The Steady State theory had trouble accounting for the helium abundance and the CMB, and was therefore abandoned for the Big Bang theory, which naturally accounted for both of these observations.

Shortly thereafter---in 1967,---a cosmological theory involving $\Lambda$ was again proposed; this time, in order to account for the apparent `preponderance' of quasars at redshift $z\sim1.95$ \cite{Petrosian1967,Shklovsky1967,Kardashev1967}. This theory---call it PS$_3$K for its authors---employed Lema\^{i}tre's cosmological model, with positive curvature and a cosmological constant, $\Lambda=\Lambda_E(1+\varepsilon)$,  $0<\varepsilon\ll1$, which described a period of stagnation during which the closed universe would slowly approach the Einstein radius, remaining there (roughly) for an arbitrary length of time---depending on how small one chose $\varepsilon$---before continuing its expansion at an accelerating rate. During this stagnation, light previously emitted by quasars would not be further redshifted, but, as it travelled through quasi-static space, would develop an absorption spectrum after passing through galaxies and interstellar dust. Thus, the theory hypothesised that the stagnation period occurred at a redshift of $z=1.95$, and sought substantiation from an observation by Geoffrey and Margaret Burbidge \cite{Burbidge1967a,Burbidge1967b}, that some quasars appear to have emission redshifts which exceed their absorption redshifts of $z\sim1.95$. However, the PS$_3$K theory never gained wide acceptance, and was completely rejected as more data and a theoretical understanding of the evolution of quasars were brought forth in the early 1970s. But it is historically significant because of its link to the first second-order calculation of the quantum vacuum energy, published later that year, by Yakov Zel'dovich \cite{Zeldovich1967,Zeldovich1968,Zeldovich2008}.

Zel'dovich \cite{Zeldovich1968,Zeldovich2008} used arguments similar to some of those we have already considered, in arguing against the Einstein interpretation of $\Lambda$; i.e., he argued that there is no theoretical justification for assuming $\Lambda\equiv0$;---that appealing to aesthetics and logical economy is a weak form of argument which could never stand up against objective data. He then pointed out, that observations of the cosmos, being always subject to the limits of experimental uncertainty, can never prove $\Lambda$ is absolutely zero, but could only ever arrive at this result gradually. Here, he made his famous statement regarding the cosmological constant: `The genie has been let out of the bottle, and it is no longer easy to force it back in.' Historically, this statement would hold true following Zel'dovich's work---for although the PS$_3$K theory would die, it had done its part in motivating Zel'dovich's calculation, which would eventually find itself as one of the pioneering works on one of the most important problems in modern physics---the cosmological constant problem; see, e.g., \cite{Rugh2002,Padmanabhan2003} for reviews of the problem from this decade. 

Interest in the vacuum energy and its possible connection to cosmology grew steadily during the 1970s, following the calculations by Zel'dovich. In particular, symmetry breaking after the hot Big Bang became a popular topic. A review of the early research in this field up to 1975 is provided in \textit{Relativistic Astrophysics, Vol.~2}, by Zel'dovich and Igor Novikov \cite{Zeldovich1983}. It describes theories which had emerged suggesting that the vacuum energy might be very large shortly after the Big Bang, due to broken symmetries, in which a link to the cosmological constant had been noted. 

Then, while working on the problem of a grand unified field theory in 1980, Alan Guth noticed that some scalar fields might get caught in a local minimum and drive an exponential inflation, from a vacuum energy that would act like a cosmological constant \cite{Weinberg2008} (which references \cite{Guth1997}). The inflation scenario, which Guth proposed the following January \cite{Guth1981}, had many problems, as he noted, and has been revised over the years; a recent description is given in \cite{Weinberg2008}. The importance of inflation theory has however remained unchanged: it proposed a solution to the horizon and flatness problems, which had been relatively unknown, but have since gained notoriety as significant problems of cosmology. 

Both problems were discussed in detail by Guth in \cite{Guth1981}. Briefly, the horizon problem is, that in the FLRW models alone, in which the cosmological spacetime geometry would be determined only by the distribution of matter, there is no way to account for the observed isotropy in the present Universe. This is because our particle horizon, or spherical region of influence, from which we currently receive light emitted at some time in the distant past, has increased linearly with time, while the FLRW scale factor has gone more like $t^{1/2}$ or $t^{2/3}$, so long as $\Lambda$ was negligible. Thus, if the expansion of space is to be determined locally by the distribution of matter, as the FLRW models assume, regions we now observe would not have been causally connected in the past, and the present distribution should not be as smoothed out as it appears. The flatness problem is that the Universe appears to have approximately zero curvature. However, according to Friedman's equations, for the Universe to be as flat as it appears today, it would've had to have been unimaginably near to being perfectly flat at very early times, as even very small deviations from flatness would lead to a present Universe which is nowhere near flat. Of course, the flatness problem vanishes if the Universe is exactly flat; but as Guth argued, `the universe is certainly not described \textit{exactly} by the Robertson-Walker metric. Thus it is difficult to imagine any physical principle which would require a parameter of that metric to be exactly equal to zero' \cite{Guth1981}.

Inflation would solve the horizon problem if a rapid expansion would have increased the size of our prior causally connected region faster than the speed of light, so that it is still much larger than our present particle horizon. This rapid expansion of space is also supposed to have flattened the curvature, like being on the surface of a sphere of increasing radius; but, as Roger Penrose has pointed out, if the matter distribution were, e.g., fractal, inflation would have had no such effect \cite{Penrose2005}. A very readable critique of the current status of the inflation scenario was recently given by Paul Steinhardt \cite{Steinhardt2011}.

Aside from the immediate significance of inflation theory in potentially solving both of these problems, it also had a hidden significance in that it was the first prediction, since the combination of the PS$_3$K theory and Zel'dovich's work, of a solution to a problem in cosmology that would be given by a theoretical calculation of vacuum energy. However, whereas the PS$_3$K theory was only short lived, the inflation theory has endured in cosmology, and has thus been a long-standing connection between cosmology and particle physics.

Throughout the 1980s, further analytical work on the vacuum energy from symmetry breakings and experimental verification of the Standard Model of particle physics---e.g.~the discovery of W and Z gauge bosons in 1983, which confirmed the electroweak interaction---led to a solidification of the theory, and evidenced a significant discrepancy between all theoretical predictions and the observational limit on $\Lambda$ from cosmology. In 1989, Steven Weinberg wrote a review \cite{Weinberg1989}, outlining the problem in which the limit on $\Lambda$ was thought to be different from the vacuum energy predictions by anywhere from forty to 120 orders of magnitude, as well as considering its various potential solutions---none of which would prove very promising. And aside from the confirmation of a small cosmological constant in 1998 \cite{Riess1998,Perlmutter1999}, which is roughly equal to the upper bound that had been known from the observed cosmic expansion, little has changed in the general picture, and no well-motivated theory has emerged with the power to properly explain this discrepancy. In 1983, Novikov asked the question, `Is the vacuum gravitating?' \cite{Novikov1983}. Due to the disparity which has only become more sure in the intervening years, it is now more common to hear the phrase, `The vacuum is gravitating extremely weakly.'

A noteworthy obstacle of the problem, is that, according to the duality of Lema\^{i}tre's interpretation of $\Lambda$, the observed dark energy can only be known as the addition of a pure $\Lambda$ and vacuum energy. For consider the vacuum Einstein equation with non-zero $\Lambda$ and with some amount of $\Lambda$-equivalent vacuum energy:
\begin{equation}
G_{\mu\nu}+g_{\mu\nu}\Lambda=\kappa{T}_{\mu\nu,\mathrm{vac}},\label{EinsteinLambdaDE}
\end{equation}
where $\kappa{T}_{\mu\nu,\mathrm{vac}}\equiv-g_{\mu\nu}\Lambda_{\mathrm{vac}}$ explicitly defines the vacuum energy as a constant source. Note, that the vacuum energy is a perfect fluid, with $T^{\mu\nu}_{\mathrm{vac}}=(\rho_{\mathrm{vac}}+p_{\mathrm{vac}})u^{\mu}u^{\nu}+ p_{\mathrm{vac}}g^{\mu\nu}$ and equation of state, $p_{\mathrm{vac}}=-\rho_{\mathrm{vac}}\equiv-\Lambda_{\mathrm{vac}}/\kappa$. Now, we are free to move $\Lambda$ to the right-hand side of Eq. (\ref{EinsteinLambdaDE}) and combine constants, defining an effective cosmological dark energy term, $\kappa{T}_{\mu\nu,\mathrm{DE}}\equiv-g_{\mu\nu}\Lambda_{\mathrm{vac}}-g_{\mu\nu}\Lambda$, with Einstein equation
\begin{equation}
G_{\mu\nu}=\kappa{T}_{\mu\nu,\mathrm{DE}}.\label{EinsteinTDE}
\end{equation}
Therefore, regardless of all possible complexities of a potentially non-zero vacuum energy, $\Lambda$ may always be absorbed into the stress energy tensor as a source. Then, if the vacuum energy is a pure constant, the observed expansion rate and our knowledge of quantum theory tell us the term $|\kappa\rho_{\mathrm{vac}}+\Lambda|$ must cancel to at least forty decimal places.

However, according to quintessence theories, this does not have to be so; for, once we have made this identification, we are free to ask whether the dark energy density really is a constant, or whether it varies slowly at sufficiently late times, so that, in the observable past, it has resembled a cosmological constant, as measurements constrain. Then, even if $\Lambda$ has a non-zero value, it could act as a potential minimum of the quintessence field. In \cite{Weinberg1989}, Weinberg notes that in order for the vacuum energy in the electroweak and earlier transitions to produce a small effective cosmological constant at late times, it is necessary to choose the potential so that earlier, before the phase transitions, the dark energy would have had a very large value, in connection with inflation theory. He has remained a strong advocate of this theory, and, in \cite{Weinberg2008}, presents an extensive analysis of quintessence models, which he feels is the surest approach to solving the cosmological constant problem.

Let us now turn to the SN Ia observations, which provide the most direct evidence of the nature of cosmic expansion: SNe Ia are explosions of white dwarf stars in binary systems following sufficient accretion of mass from their companions to reach the limiting mass that can be supported by degenerate electron pressure, known as the Chandrasekhar limit \cite{Chandrasekhar1931a,Chandrasekhar1935}, so that their constituent carbon would necessarily ignite explosively. Thus, the luminosity curves of these explosions are always the same, aside from measurable perturbances, and they act as `standard candles'; i.e., when their apparent brightnesses have been measured, their distances can be calculated quite accurately. By observing numerous SNe Ia at various redshifts, astronomers have therefore been able to fairly accurately trace the changing rate of cosmic expansion, and the result, which has been known for more than a decade now \cite{Riess1998,Perlmutter1999}, is that the cosmic expansion is accelerating.

From Eqs. (\ref{Friedmann}) -- (\ref{Friedmann2}), we have the acceleration of expansion in the FLRW universes,
\begin{equation}
\frac{\ddot{a}}{a}=\frac{\Lambda}{3}-\frac{\kappa}{2}\left(p+\frac{\rho}{3}\right).\label{ddot_a}
\end{equation}
Therefore, in any FLRW model, a measurement of positive acceleration requires positive $\Lambda$, in order to eventually overcome the decreasing contribution to the expansion rate from matter---for which $\rho,~p\geq0$---at late times. In 2007, the ESSENCE (`Equation of State: SupErNovae trace Cosmic Expansion') SN~Ia data were found to place significant statistical constraint on the restriction of exotic dark energy models to the so-called flat $\Lambda$CDM model, which describes a flat FLRW universe dominated by a pure cosmological constant and cold dark matter \cite{Wood-Vassey2007,Davis2007}. The constraint on dark energy was strengthened in subsequent years, with the publications of results evaluated using the Union \cite{Kowalski2008} and Constitution \cite{Hicken2009} SN~Ia data sets. More recently, a claim was presented in the publications of cosmological results from the first year Sloan Digital Sky Survey-II (SDSS-II) SN data release \cite{Kessler2009,Sollerman2009}, that these data may favour more exotic dark energy; however, there is evidence that this result is artificial \cite{Bueno2009}.

Measurements of anisotropies in the CMB and baryon acoustic oscillations (BAO) together provide constraints on cosmological parameters which are independent from those found by analysis of SN Ia data, and which have also always been consistent with the flat $\Lambda$CDM model; and, as with the SN Ia data, the most recent results from the Wilkinson Microwave Anisotropy Probe (WMAP) 7-year data \cite{Jarosik2011,Komatsu2011} and the spectroscopic SDSS Data Release 7 galaxy sample \cite{Percival2010} have significantly constrained those parameters.

Thus, the current status of the cosmological constant problem is that quantum theory predicts a vacuum energy which is much larger than the observed value of $\Lambda_{\mathrm{DE}}$, and there is yet no convincing evidence that a decay to that value should have occurred. Furthermore, the value of $\Lambda_{\mathrm{DE}}$ measured from SN~Ia light curves, as well as from CMB and BAO measurements, has the same order of magnitude as the indirect constraint which had previously been placed on it through expansion dynamics; viz., $\Lambda_{\mathrm{DE}}\approx10^{-52}~\mathrm{m}^{-2}$. This constitutes half of the `coincidence problem':---that the matter and vacuum densities are roughly equal at present. For if $\Lambda_{\mathrm{DE}}$ were truly unrelated to the primordial cosmic expansion, it might instead have been that $H_0\gg\left|\sqrt{\Lambda_{\mathrm{DE}}}\right|$, rather than the two terms being nearly equal; then, even with the precision measurements from SNe~Ia and the CMB, $\Lambda_{\mathrm{DE}}$ would have gone undetected. On the other hand, $\Lambda_{\mathrm{DE}}$ might have been very large, so that the Universe would have quickly become de~Sitter's vacuum, with matter density negligible in comparison, and the Universe would now appear much different to us; e.g., if galaxies \textit{could} have formed, they may all have vanished from sight long ago, and we would never have detected the Universal expansion. Such possibilities motivate anthropic arguments, as, e.g., Weinberg has considered \cite{Weinberg1989}.

Altogether, these results suggest that the measured value of $\Lambda_{\mathrm{DE}}$ is not likely related to the predicted vacuum energy from quantum field theory, but that our Universe has a true cosmological constant---which could rather be essentially linked to its apparent expansion. If this were true, the implication could be that we must rectify quantum field theory with vanishing vacuum energy, or else our physical cosmology would have to require that the large vacuum energy does not contribute to the cosmic expansion.

Now, the principal issue with any attempt to reconcile Eddington's interpretation of expansion with theory, is that it is inconsistent with FLRW big bang cosmology, according to Eq.~(\ref{ddot_a}). As we shall eventually see, however, the RW line-element corresponds only the most trivial first guess that may be made in order to formulate a dynamical model to account for our cosmological observations. And this first guess, which should anyhow be useful for modelling those \textit{observations} from anywhere in the Universe, according to our own observations and the cosmological principle, has since been investigated extensively, with the uncomfortable result, that when we further interpret the observational model as a representation of the Real physical state of the Universe, the theory is plagued with many significant problems to which we have been unable to find truly sufficient solutions. 

Most notable among these, is the frequently overlooked problem that was known to Eddington:---that, in the big bang scenario, FLRW cosmology cannot \textit{explain} the expansion we use the model to empirically \textit{describe}; for only Weyl's `de~Sitter cosmology', completely devoid of matter, is capable of explaining the expansion from the moment of a big bang through a positive cosmological constant, while all other FLRW big bang universes require an indescribable singular push in order to exist. Thus, these models can describe only the effects of extrinsic sources on a primordial expansion rate, and provide no realistic clues as to its origin.

The Nobel Prize-winning discovery late last century, that the cosmic expansion is accelerating, was therefore a nuisance for standard cosmology because it adds a complication to the logical simplicity of the theory; viz., that the measured positive $\Lambda_{\mathrm{DE}}$ has no fundamental role in the theory, being unnecessary for the existence of a FLRW big bang universe, and so begs the question, `Why should it exist at all?' $\Lambda_{\mathrm{DE}}$ has therefore been popularly thought of as vacuum energy, which could have had a role in the inflation scenario if the quintessence theory is correct. Otherwise, only at late times would pure $\Lambda$ become relevant for describing cosmic expansion. 

In essence, in the standard theory, the observation of cosmic acceleration is an annoyance, and the measured value of $\Lambda_{\mathrm{DE}}$ is a problem. And the logical simplicity that Einstein had sought while supporting the Einstein-de~Sitter universe was in fact only superficial, as it imposes on cosmology an impossible complication, in which the entire history of the Universe is ultimately to be the result of an inexplicable singularity, at which point (we are then forced to conclude), the theory itself must fail to describe things correctly; thus, Eddington's interpretation is logically much simpler because, at the very least, it does not need to presuppose that the mathematical theory needs to break down at the single event that is also supposed to essentially give rise to our extraordinarily particular reality (for more discussion on the extraordinary specialness of the particular big bang that this theory should require, see \cite{Penrose2005,Penrose2009}).

It is no wonder that Eddington was so incensed with the position Einstein had finally taken regarding the observed expansion, which came only after Hubble's confirmation, after he had remained silent on the issue during the prior decade's theoretical investigation by so many others: the feeling surely would have been similar for Gamow, Alpher, Herman, and Lema\^{i}tre, if, in 1965, Robert Dicke and his colleagues had, instead of \cite{Dicke1965}, written that Penzias and Wilson's discovery confirmed a prediction of the Steady State theory, in the particular case in which God continually fills the Universe with a steady amount of one-quarter helium, three-quarters hydrogen, and three-degrees of microwave radiation; for the \textit{explanation} of expansion, according to a prior physical cause, was similarly unjustified in the Einstein-de~Sitter model, which not only completely neglected the manner in which the observed expansion had been \textit{anticipated} through physical theory, but even neglected any mention of the possibility.

Einstein seems to have had little interest in understanding prior physical causes, as a means of developing a general relativistic cosmological theory that would reasonably account for the observed phenomena: instead, he seems to have been primarily concerned with the construction of an empirical model that would fit the observations in a manner that would remain consistent with his previous reasoning---i.e., he really didn't give much further thought to the problem of expansion. His evident lack of interest in finding a prior explanation to this problem is indeed a concern; and the best explanation that I can see for this apparent apathy, is that he had already led himself towards more abstract ways of thinking, due to the idealist-determinist interpretation of the meaning of relativity that he always considered to be the most objective one, which led him to accept, sometime long before the 1930s---e.g., as already evidenced in the postcard he sent Weyl in May, 1923,---that the general theory must be essentially flawed.

The abstraction of thought to which I am referring is evident, e.g., in his autobiography, where he wrote that he spent seven difficult years (1908-1915) freeing himself `from the idea that co-ordinates must have an immediate metrical meaning' \cite{Einstein1951}. However, in contrast to the way of conceiving things he thus came to, the coordinate systems used to describe physical reality \textit{do} have a clear metrical meaning in the dynamical theory of relativity proposed in Chapter~\ref{chap_DIRT}, which describes the relativistic spacetime of empirical observation as an ideal graph of events which occur in an objectively well-defined three-dimensional universe that physically evolves in an objectively well-defined manner.

But so our investigation must first turn to the problem of understanding why Einstein had denied this possibility, and accepted instead an abstract interpretation of the relativistic description of events in physical reality that ultimately requires four-dimensional block determinism; and we shall see that he was in fact entirely consistent in that interpretation---and, therefore, that the common application of relativity theory, which has always demanded the same interpretation of the mathematics that \textit{does} formally require spacetime to be a singular block eternity, has not been so consistent. Accordingly, then, it shall become clear that even though relativity theory is commonly used to describe dynamic processes, the inherent inconsistencies in such applications have obstructed an accurate understanding of the true dynamical description that is afforded by the theory, and many paradoxes have subsequently been inferred which we shall immediately see resolved when the formally consistent dynamical interpretation of the mathematical theory is assumed.

So, in order to completely justify that interpretation as the most objective one, we must thoroughly assess the many aspects of this basic problem of four-dimensionalism, and try to understand the assumptions and presumptions behind the reasonings of those who have played important roles in its philosophical development, with the idea that if we can truly understand what they were thinking and why, we might better understand where they went wrong. For only then, when we have subsequently applied that knowledge in consistently working through the formal requirements for causal dynamism, will we be able to properly address the cosmological problem of relativity theory in a straightforward, logically consistent manner, and work out a solution that is consistent with cosmological observations, the mathematical theory, and the true dynamical nature of reality we want our theory to describe---which, we'll find, objectively accounts for the details demanded by both reason and observation.